\documentclass[11pt]{article}

% ===================== PACKAGES =====================
\usepackage[a4paper,margin=1in]{geometry}
\usepackage{amsmath,amssymb}
\usepackage{enumitem}
\usepackage{fancyhdr}

% ===================== HEADER =====================
\pagestyle{fancy}
\fancyhf{}
\rhead{CSE400}
\lhead{Lecture 19}
\rfoot{\thepage}

\begin{document}

\begin{center}
    {\LARGE \textbf{Lecture 19: N-Variate Joint PDF/PMF/CDF}}\\[0.5em]
    {\large CSE400 - Fundamentals of Probability in Computing}\\[0.5em]
    {\normalsize Dhaval Patel, PhD}\\
    {\normalsize SEAS, Ahmedabad University}\\
    {\normalsize March 17, 2026}
\end{center}

\vspace{1em}

% ===================== OUTLINE =====================
\section*{Outline}
\begin{itemize}
    \item Joint Distribution of $N$ Random Variables
    \begin{itemize}
        \item Definition
        \item Example-1
        \item Application: Large MIMO System
        \item Joint PMF, CDF, PDF
        \item Independence of Random Variables
        \item Example-2: Uniformly Distributed 3D Vector
        \item Example-3: Trivariate Random Variables
    \end{itemize}
\end{itemize}

% ===================== SECTION =====================
\section{Joint Distribution of $N$ Random Variables}

\subsection{Definition}

Let $X_1, X_2, \dots, X_N$ be $N$ random variables defined on the same probability space.

\begin{itemize}
    \item The \textbf{joint cumulative distribution function (CDF)} is defined as:
    \[
    F_{X_1, X_2, \dots, X_N}(x_1, x_2, \dots, x_N)
    = P(X_1 \le x_1, X_2 \le x_2, \dots, X_N \le x_N)
    \]
\end{itemize}

\subsection{Joint PMF}

For discrete random variables:

\begin{itemize}
    \item The \textbf{joint probability mass function (PMF)} is:
    \[
    p_{X_1, X_2, \dots, X_N}(x_1, x_2, \dots, x_N)
    = P(X_1 = x_1, X_2 = x_2, \dots, X_N = x_N)
    \]
\end{itemize}

\subsection{Joint PDF}

For continuous random variables:

\begin{itemize}
    \item The \textbf{joint probability density function (PDF)} is defined such that:
    \[
    F_{X_1, \dots, X_N}(x_1, \dots, x_N)
    = \int_{-\infty}^{x_1} \cdots \int_{-\infty}^{x_N}
    f_{X_1, \dots, X_N}(t_1, \dots, t_N)\, dt_1 \cdots dt_N
    \]
\end{itemize}

\subsection{Properties}

\begin{itemize}
    \item Non-negativity:
    \[
    f_{X_1, \dots, X_N}(x_1, \dots, x_N) \ge 0
    \]
    
    \item Total probability:
    \[
    \int \cdots \int f_{X_1, \dots, X_N}(x_1, \dots, x_N)\, dx_1 \cdots dx_N = 1
    \]
\end{itemize}

% ===================== EXAMPLE 1 =====================
\section{Example-1}

(As presented in slides)

\begin{itemize}
    \item Consider joint behavior of multiple random variables.
    \item Evaluate probabilities using joint PMF/PDF definitions.
\end{itemize}

% ===================== APPLICATION =====================
\section{Application: Large MIMO System}

\begin{itemize}
    \item Multiple-input multiple-output (MIMO) systems involve multiple random variables.
    \item Signal components across antennas can be modeled as joint random variables.
    \item Joint distributions are required to analyze system performance.
\end{itemize}

% ===================== INDEPENDENCE =====================
\section{Independence of Random Variables}

\subsection{Definition}

Random variables $X_1, X_2, \dots, X_N$ are independent if:

\[
f_{X_1, X_2, \dots, X_N}(x_1, x_2, \dots, x_N)
= \prod_{i=1}^{N} f_{X_i}(x_i)
\]

for continuous case, or

\[
p_{X_1, X_2, \dots, X_N}(x_1, x_2, \dots, x_N)
= \prod_{i=1}^{N} p_{X_i}(x_i)
\]

for discrete case.

\subsection{Equivalent Condition using CDF}

\[
F_{X_1, X_2, \dots, X_N}(x_1, \dots, x_N)
= \prod_{i=1}^{N} F_{X_i}(x_i)
\]

% ===================== EXAMPLE 2 =====================
\section{Example-2: Uniformly Distributed 3D Vector}

\begin{itemize}
    \item Consider a 3D random vector:
    \[
    (X, Y, Z)
    \]
    
    \item Uniformly distributed over a region in 3D space.
    
    \item Joint PDF:
    \[
    f_{X,Y,Z}(x,y,z) = \text{constant}
    \]
    over the valid region.
    
    \item Outside the region:
    \[
    f_{X,Y,Z}(x,y,z) = 0
    \]
    
    \item Constant determined using normalization:
    \[
    \iiint f_{X,Y,Z}(x,y,z)\,dx\,dy\,dz = 1
    \]
\end{itemize}

% ===================== EXAMPLE 3 =====================
\section{Example-3: Trivariate Random Variables}

\begin{itemize}
    \item Consider three random variables:
    \[
    X, Y, Z
    \]
    
    \item Joint distribution defined using:
    \[
    f_{X,Y,Z}(x,y,z)
    \]
    
    \item Marginal distributions obtained by integration:
    \[
    f_X(x) = \int \int f_{X,Y,Z}(x,y,z)\, dy\, dz
    \]
    
    \[
    f_Y(y) = \int \int f_{X,Y,Z}(x,y,z)\, dx\, dz
    \]
    
    \[
    f_Z(z) = \int \int f_{X,Y,Z}(x,y,z)\, dx\, dy
    \]
\end{itemize}

\subsection{General Marginalization}

\begin{itemize}
    \item For any subset of variables:
    \[
    f_{X_1, \dots, X_k}(x_1, \dots, x_k)
    = \int \cdots \int f_{X_1, \dots, X_N}(x_1, \dots, x_N)\, dx_{k+1}\cdots dx_N
    \]
\end{itemize}

% ===================== END =====================
\end{document}